\documentclass[10pt,twocolumn]{article}

\usepackage[utf8]{inputenc}
\usepackage[T1]{fontenc}
\usepackage{lmodern}
\usepackage[margin=0.75in,top=1in,bottom=1in]{geometry}
\usepackage{amsmath,amssymb}
\usepackage{booktabs}
\usepackage{graphicx}
\usepackage{xcolor}
\usepackage{hyperref}
\usepackage{microtype}
\hypersetup{colorlinks=true,linkcolor=blue!60!black,citecolor=green!50!black,urlcolor=blue!60!black}

\usepackage{cerebrum-macros}
% Unified notation table for CEREBRUM papers
% Resolve naming conflicts: TSC uses \eta_T for temporal weight; CSA uses \eta for the same concept.
% All papers should import this file and use the macros below for consistency.

\newcommand{\etaT}{\eta_T}        % TSC temporal weight (renamed from generic \eta to avoid conflict)
\newcommand{\etaCSA}{\eta}        % CSA temporal decay weight (10-param formula, position 7)
\newcommand{\alphaCSA}{\alpha}    % CSA semantic similarity weight
\newcommand{\betaCSA}{\beta}      % CSA community score weight
\newcommand{\gammaCSA}{\gamma}    % CSA edge-type weight
\newcommand{\deltaCSA}{\delta}    % CSA distance penalty
\newcommand{\epsilonCSA}{\epsilon}% CSA hop decay
\newcommand{\zetaCSA}{\zeta}      % CSA PageRank prior weight
\newcommand{\iotaCSA}{\iota}      % CSA node recency weight
\newcommand{\muCSA}{\mu}          % CSA synthesis-density penalty
\newcommand{\thetaCSA}{\theta}    % CSA grounding confidence weight

% Graph notation
\newcommand{\KG}{\mathcal{G}}     % Knowledge graph
\newcommand{\Eset}{\mathcal{E}}   % Edge set
\newcommand{\Vset}{\mathcal{V}}   % Vertex/node set
\newcommand{\Cset}{\mathcal{C}}   % Community set
\newcommand{\csascore}{a(u,v,k)}  % CSA attention score

% Traversal notation
\newcommand{\beam}{B}             % Beam width
\newcommand{\hops}{H}             % Number of hops
\newcommand{\mrr}{\text{MRR}}     % Mean Reciprocal Rank
\newcommand{\hitatk}[1]{\text{H@#1}} % Hits@k


\title{\textbf{Holographic Indexing: Privacy-Preserving Discovery\\
in Federated Knowledge Graph Networks}}

% Author block — shared across all CEREBRUM arXiv papers
\author{
  Bryan Alexander Buchorn \\
  Independent Researcher \\
  Las Vegas, NV, USA \\
  \texttt{bryan.buchorn@gmail.com}
}

\date{May 2026}

\begin{document}
\maketitle

% -----------------------------------------------------------------------
\begin{abstract}
Federated Knowledge Graph reasoning requires nodes to identify relevant peers without
revealing the contents of their local graphs. We present \textbf{Holographic Indexing},
a two-tier privacy-preserving discovery protocol combining (1) Bloom Filter membership
probing for exact entity matching and (2) Community Centroid Signatures for semantic
neighborhood approximation. The ``holographic'' representation allows a querying node
to identify potential reasoning destinations in remote peers with bounded information
leakage: exact entity membership cannot be enumerated from a Bloom filter, and centroid
signatures reveal only community-level topology, not individual entities. We formalize
the privacy properties (Theorem~\ref{thm:privacy}), analyze communication complexity
($O(|C| \cdot d + m)$ bits per peer probe vs. $O(|\Vset|)$ for full entity exchange),
and describe Orthogonal Procrustes alignment for cross-node embedding space
normalization. The protocol is implemented in the \CEREBRUM{} framework via
\texttt{DistributedBeamTraversal}, with HMAC-SHA256 path provenance ensuring integrity
of federated reasoning paths. We further show that the same discovery layer generalizes
to scientific literature federation via the \texttt{ExternalValidator} component.
\end{abstract}

% -----------------------------------------------------------------------
\section{Introduction}
\label{sec:intro}

Federated Knowledge Graphs are increasingly common: healthcare systems partition patient
data across hospital nodes; enterprise knowledge bases span business units; scientific
knowledge is distributed across domain-specific repositories. Cross-node reasoning
requires identifying which peers hold relevant information — but disclosing full entity
lists violates data sovereignty constraints and creates enumeration attack surfaces.

Traditional federated search approaches require either a central index (single point of
failure, privacy violation) or full entity exchange (prohibitive bandwidth, privacy
violation). We propose a middle path: a \emph{holographic} representation that enables
semantic discovery without enabling entity enumeration.

The term ``holographic'' captures the key property: like a holographic projection, the
representation contains enough information to identify relevant regions, but destroying
the projection does not recreate the full original. A Bloom filter over entity IDs
cannot be inverted; a community centroid cannot recreate individual node embeddings.

\paragraph{Contributions.}
\begin{enumerate}
  \item Tier-1 Bloom Filter probing: exact membership testing with configurable
        false-positive rate $p$ and no enumeration leakage.
  \item Tier-2 Community Centroid Signatures: semantic neighborhood approximation
        via cosine similarity to centroid/radius pairs.
  \item Formal privacy analysis (Theorem~\ref{thm:privacy}): bounded information
        leakage for both tiers.
  \item Communication complexity: $O(|C| \cdot d + m)$ vs $O(|\Vset|)$ per query.
  \item Orthogonal Procrustes alignment for cross-node embedding normalization.
  \item HMAC-SHA256 path provenance preventing adversarial path injection.
\end{enumerate}

% -----------------------------------------------------------------------
\section{Background}
\label{sec:background}

\paragraph{Bloom filters.}
A Bloom filter~\cite{bloom1970} is a space-efficient probabilistic data structure for
set membership testing. An element is inserted by hashing it with $k$ hash functions
and setting corresponding bits. Membership is tested by checking all $k$ bit positions.
False positives occur with rate $p \approx (1-e^{-kn/m})^k$ for $n$ elements in $m$
bits; false negatives are impossible. Crucially, the full set $\mathcal{E}$ cannot be
enumerated from the filter.

\paragraph{Federated learning and KG federation.}
\citet{kairouz2021} survey open problems in federated learning; \citet{bonatti2019} address
federated SPARQL query answering. Existing KG federation protocols (OWL-based,
SPARQL endpoints) typically expose full query capability, not bounded-leakage discovery.

\paragraph{Embedding alignment.}
When two nodes train embeddings independently, their spaces are not comparable even for
shared entities. The Orthogonal Procrustes problem~\cite{schonemann1966} finds the
rotation $R^* = \arg\min_R \|XR - Y\|_F$ s.t. $R^TR = I$, solved analytically via SVD.

% -----------------------------------------------------------------------
\section{Holographic Indexing Protocol}
\label{sec:protocol}

\subsection{Tier 1: Entity Bloom Filter}

Each node $\mathcal{G}_i$ constructs a Bloom filter $B_i$ over its entity set $\mathcal{E}_i$:
\begin{equation}
B_i = \mathrm{BloomFilter}(\mathcal{E}_i,\; m_i,\; k)
\label{eq:bloom}
\end{equation}

A querying node $\mathcal{G}_j$ probes $B_i$ for a specific entity $e$:
\begin{equation}
\Pr[\mathrm{FP}] = \left(1 - e^{-k\,|\mathcal{E}_i|/m_i}\right)^k
\label{eq:fpr}
\end{equation}

Default parameters: $m_i = 10 \cdot |\mathcal{E}_i|$ bits, $k = 7$ hash functions, yielding
$\Pr[\mathrm{FP}] \approx 0.008$. Bloom filter size is $O(|\mathcal{E}_i|)$ bits,
transmitted once per session and cached locally.

\subsection{Tier 2: Community Centroid Signatures}

For semantic discovery (find peers with \emph{similar} entities), entity-level probing
is insufficient. We introduce Community Centroid Signatures: for each community
$C_k \in \mathcal{G}_i$:
\begin{equation}
\text{sig}(C_k) = (\vec{c}_k,\; \rho_k) \quad\text{where}\quad
\vec{c}_k = \mathrm{mean}\{\vec{e} : e \in C_k\},\;\;
\rho_k = \max_{e \in C_k} \|\vec{e} - \vec{c}_k\|
\label{eq:signature}
\end{equation}

A querying node with entity embedding $\vec{x}$ computes the \textbf{Synaptic Bridge Score}:
\begin{equation}
\mathrm{SBS}(C_k) = \cos(\vec{x}, \vec{c}_k)
\label{eq:sbs}
\end{equation}

Peers with $\mathrm{SBS}(C_k) \geq \sigma_{thresh}$ (default 0.75) are flagged as
relevant reasoning destinations. Only $|C_i| \cdot (d + 1)$ floats are transmitted,
where $|C_i|$ is the number of communities and $d$ is the embedding dimension.

\subsection{Communication Complexity}

\begin{table}[h]
\small\centering
\caption{Communication cost per query per peer.}
\label{tab:comm}
\begin{tabular}{lc}
\toprule
\textbf{Method} & \textbf{Bits transmitted} \\
\midrule
Full entity list       & $O(|\Vset| \cdot \log|\Vset|)$ \\
Full entity embeddings & $O(|\Vset| \cdot d)$ \\
\midrule
Tier 1 (Bloom)         & $O(|\Vset|)$ (one-time) \\
Tier 2 (centroids)     & $O(|C| \cdot d)$ (one-time) \\
Query (both tiers)     & $O(1)$ probe + $O(|C_{hit}|)$ path delegation \\
\bottomrule
\end{tabular}
\end{table}

After the initial handshake, per-query communication is $O(1)$ for discovery and
$O(|C_{hit}|)$ for reasoning — orders of magnitude less than full entity exchange.

% -----------------------------------------------------------------------
\section{Privacy Analysis}
\label{sec:privacy}

\begin{theorem}[Bounded Leakage]
\label{thm:privacy}
The Holographic Index $\mathcal{H}_i = (B_i, \{\mathrm{sig}(C_k)\}_k)$ satisfies:
\begin{enumerate}
  \item \textbf{Enumeration resistance}: $\mathcal{E}_i$ cannot be recovered from $B_i$
        in polynomial time (follows from the pre-image resistance of the hash functions).
  \item \textbf{Bounded semantic leakage}: An adversary can identify communities whose
        centroid is within angle $\theta$ of a target vector, but cannot identify
        individual entities within those communities.
  \item \textbf{False positive cover}: Any entity $e \notin \mathcal{E}_i$ has probability
        $p < 0.01$ of appearing to be present, providing plausible deniability for
        boundary cases.
\end{enumerate}
\end{theorem}

\textit{Proof.} (1) Inverting a Bloom filter requires $O(2^m)$ preimage queries;
computationally infeasible for $m > 256$ bits. (2) Centroid signatures are one-way
aggregations: a centroid $\vec{c}_k = \mathrm{mean}\{\vec{e}_j\}$ has entropy
$H(\mathcal{E}_{C_k}) - H(\vec{c}_k) > 0$ for $|C_k| > 1$, i.e., individual
embeddings cannot be recovered without additional information. (3) Follows directly
from Bloom filter false-positive probability~(\ref{eq:fpr}). $\square$

\paragraph{Practical limitations.}
Centroids reveal approximate community topic. In a healthcare KG, a centroid near
``cardiology'' reveals that the node holds cardiology entities, without revealing
which patients. Whether this constitutes acceptable leakage is domain-specific.
For higher-sensitivity deployments, differentially private noise can be added to
centroids~\cite{dwork2014dp} at a cost to matching accuracy.

% -----------------------------------------------------------------------
\section{Cross-Node Embedding Alignment}
\label{sec:procrustes}

Community centroid similarity scores are only meaningful if the two nodes share a common
embedding space. When nodes train embeddings independently (or use different models),
a calibration step is required.

Given shared ``anchor'' entities $\mathcal{A} = \mathcal{E}_i \cap \mathcal{E}_j$,
we solve:
\begin{equation}
R^* = \arg\min_{R:\,R^TR=I} \|X_\mathcal{A} R - Y_\mathcal{A}\|_F
\label{eq:procrustes}
\end{equation}
where $X_\mathcal{A}, Y_\mathcal{A}$ are the anchor embeddings in nodes $i$ and $j$
respectively. The analytical solution is $R^* = VU^T$ where $X_\mathcal{A}^T Y_\mathcal{A}
= U\Sigma V^T$ (SVD). Node $i$ rotates its embeddings by $R^*$ before computing SBS scores
against node $j$'s centroids.

If the anchor set is small ($|\mathcal{A}| < 10$), alignment is unreliable. In this
case, the system falls back to Tier-1 only (exact Bloom probing) and skips semantic
matching.

% -----------------------------------------------------------------------
\section{Security: HMAC Path Provenance}
\label{sec:security}

A compromised or malicious federated node could inject high-confidence fake paths.
We prevent this with HMAC-SHA256 path provenance: each reasoning response from a
remote adapter is signed with a pre-shared key $\mathcal{K}$:
\begin{equation}
\mathrm{sig}_{path} = \mathrm{HMAC\text{-}SHA256}(\mathcal{K},\; \mathrm{path\_data})
\label{eq:hmac}
\end{equation}

The coordinator verifies the signature before merging the remote path into the beam.
Federated leases include TTL-bounded tokens ($\mathrm{exp} = t + 3600s$) to prevent
replay attacks. Paths without valid signatures are discarded — they do not degrade
reasoning quality, only coverage.

% -----------------------------------------------------------------------
\section{Scientific Literature Federation}
\label{sec:extvalidator}

The same discovery layer generalizes beyond peer KG nodes. The \texttt{ExternalValidator}
extends Holographic Indexing to open scientific literature: given a proposed hypothesis
edge $(u, \texttt{CAUSES}, v)$ from the \texttt{HypothesisEngine}, \texttt{ExternalValidator}
queries PubMed, arXiv, and OpenAlex for supporting or contradicting papers.

This represents a qualitative extension of ``federated knowledge discovery'' from
structured peer graphs to unstructured text corpora. The Bloom Filter tier is replaced
by TF-IDF entity matching; the Centroid Signature tier is replaced by embedding
similarity against paper abstract vectors. The protocol structure — two-tier probe
followed by path delegation — is preserved.

% -----------------------------------------------------------------------
\section{Conclusion}
\label{sec:conclusion}

Holographic Indexing provides a privacy-preserving, communication-efficient protocol for
federated KG discovery. The two-tier architecture (Bloom filter + community centroids)
enables semantic reasoning across organizational boundaries without entity enumeration.
Formal privacy guarantees, Procrustes embedding alignment, and HMAC path provenance make
the protocol suitable for privacy-sensitive production deployments. The generalization to
scientific literature federation demonstrates that the protocol extends naturally beyond
structured peer graphs.

\section*{Acknowledgments}
The author gratefully acknowledges the use of Claude (Anthropic) as a research assistant
for mathematical formalization, code generation, and manuscript preparation. All
conceptual contributions, architectural decisions, experimental design, and intellectual
claims are solely the author's.

\bibliographystyle{plain}
\begin{thebibliography}{99}

\bibitem{bloom1970}
Bloom, B.H. (1970).
Space/time trade-offs in hash coding with allowable errors.
\emph{Communications of the ACM}, 13(7).

\bibitem{kairouz2021}
Kairouz, P., et al. (2021).
Advances and open problems in federated learning.
\emph{Foundations and Trends in Machine Learning}, 14(1-2).

\bibitem{bonatti2019}
Bonatti, P.A., et al. (2019).
The SPECIAL-K personal data processing transparency and compliance system.
\emph{ISWC}.

\bibitem{schonemann1966}
Schönemann, P.H. (1966).
A generalized solution of the orthogonal Procrustes problem.
\emph{Psychometrika}, 31(1).

\bibitem{dwork2014dp}
Dwork, C., \& Roth, A. (2014).
The algorithmic foundations of differential privacy.
\emph{Foundations and Trends in Theoretical Computer Science}, 9(3-4).

\bibitem{broder2004}
Broder, A., \& Mitzenmacher, M. (2004).
Network applications of Bloom filters: A survey.
\emph{Internet Mathematics}, 1(4).

\bibitem{buchorn2026flagship}
Buchorn, B.A. (2026).
CEREBRUM: Training-Free Knowledge Graph Reasoning via Community-Structured Graph Attention.
arXiv preprint.

\bibitem{buchorn2026report}
Buchorn, B.A. (2026).
CEREBRUM v2.51: Complete Technical Specification.
arXiv preprint [CEREBRUM\_REPORT\_ID].

\end{thebibliography}

\end{document}
